\setvariables[meta]
  [title={Wellenoptik: Radar und Lidar},
   subtitle={Bearbeitung des Arbeitsblatt: PhQ_WO_01_RadarLidar},
   date={2024-10-20},
   author={Niklas von Hirschfeld},
  ]
\setupinteraction[title={Wellenoptik: Radar und Lidar}]

\setuphead[chapter][
	header=normal,
	page=no
	before={\testpage[12]\blank[big]},
	after={\blank[2*big]},
]

\starttext

\startfrontmatter
\component c_titlepage
\component c_toc
\stopfrontmatter

\setuppagenumber[number=1]

\startbodymatter
\chapter[title={Aufgabe 1},reference={aufgabe-1}]

\startframedtask
Recherchieren Sie RADAR und LIDAR (Quellen angeben!).

Beschreiben Sie (kurz), wie diese Technologien funktionieren und für
welche Zwecke sie eingesetzt werden.

\stopframedtask

Lidar, Radar, oder auch Sonar, funktionieren alle sehr ähnlich, indem
Energiewellen ausgesendet werden und auf eine Reflexion gewartet wird.

\section[title={Radar},reference={radar}]

\quotation{{\bf Ra}dio {\bf d}etection {\bf a}nd {\bf r}anging}, kurz
Radar, ist die Bezeichnung für verschiedene Erkennungs- oder auch
Ortungsverfahren, welche mit {\bf elektromagnetischen Wellen} im
Radiofrequenzbereich (kleiner als
$300\text{MHz}$\footnote{\cite[alternative=entry][Radiowellen_LEIFI]})

\subsection[title={Funktionsweise},reference={funktionsweise}]

Für die Funktionsweise von Radar-Sensoren sind drei grundlegende
physikalische Gesetze von
bedeutung\footnote{\cite[alternative=entry][Wolff2024Aug]}:

\startitemize[packed]
\item
  Die {\bf Reflexion} elektromagnetischer Wellen.
\item
  Die {\bf konstante Ausbreitungsgeschwindigkeit} der
  elektromagnetischen Wellen, welche sich mit knapp Lichtgeschwindigkeit
  ausbreiten (ca $300 \frac{km}{s}$).
\item
  Die {\bf geradlinige Ausbreitung} der elektromagnetischen Wellen.
\stopitemize

Die elektromagnetischen Wellen werden also ausgesendet und deren
Reflexion gemessen. Mit der Zeitdifferenz des aussenden und des
empfangen kann die Distanz ermittelt werden, welche die Wellen
zurückgelegt haben.

\subsection[title={Anwendung},reference={anwendung}]

Die Anwendungsmöglichkeiten sind vielfältig. Die Technologie wird unter
anderem dafür genutzt, um Positionen zu bestimmen, oder auch füllstände
in Lagertanks\footnote{\cite[alternative=entry][IFM_Radar]}

\section[title={Lidar},reference={lidar}]

„{\bf Li}ght {\bf d}etection {\bf a}nd {\bf r}anging“, kurz Lidar, ist
eine Technologie, welche die präziese messung von Entfernungen mit sehr
geringer latenz erlaubt\footnote{\cite[alternative=entry][IBM_LIDAR]}.

\subsection[title={Funktionsweise},reference={funktionsweise-1}]

Klassische Lidar-Sensoren bezitzen einen Transmitter und einen
Empfänger. Über den Transmitter wird ein {\bf gepulster Laser} wird
ausgesendet. Der Zeitpunkt dessen Reflexion wird mit dem Empfänger
gemessen und somit die \quotation{Flugdauer} ermittelt. Da die
Lichtgeschwindigkeit konstant ist (ca $300
\frac{km}{s}$\footnote{\cite[alternative=entry][lichtgeschwindigkeit_LEIFI]})
ist diese Methode sehr schnell und präziese.

\

\subsection[title={Anwendung},reference={anwendung-1}]

Lidar Technologie findet in der heutigen Welt eine breite und
vielschichtige Anwendung. Von Forschung bis Technologie im Alltag,
Lidar-Sensoren sind überall zu finden und auch nicht mehr weg zu denken.

In der Archäologie wurde 2018 mithilfe von Lidar-Luftscans eine bisher
übersehene Maya-Ruinenstadt gefunden. Dank der Scans konnte eine
großflächige und detailreiche 3D-Karte kreiert werden, welche die
Umrisse der Ruinen
aufdeckte\footnote{\cite[alternative=entry][lidar_defunk]}.

Auch für das autonome Fahren spielt die Lidar-Technologie eine große
Rolle. Die Autos müssen ihre Umgebung in drei Dimensionen erfassen
können. Die Lidar-Sensoren werden mit vielen anderen Sensoren, wie auch
Kameras, kombiniert um die aktuelle Situation einzuschätzen. Lidar ist
deshalb geeignet, da es sehr präzise Messungen mit sehr geringer Latenz
ermöglicht\footnote{\cite[alternative=entry][9127855]}. Aber nicht nur
für das Autonome steuern von Autos wird diese Technologie genutzt, auch
Staubsaug-Roboter nutzen sie, um sich in einem Raum zurecht zu finden.


\chapter[title={Aufgabe 2},reference={aufgabe-2}]

\startframedtask
Erklären Sie die Funktionsweise dieser beiden Technologien physikalisch.
Geben Sie dafür auch die Gleichungen an, mit denen die nötigen
Berechnungen durchgeführt werden können. Verzichten Sie auf Erklärungen
zur Nutzung der Raman-Streuung mit LIDAR-Systemen

\stopframedtask

Die Messung der \quotation{Time of flight}, also der Strecke, die die
Energiewelle in welcher Zeit zurückgelegt hat, sieht wie folgt
aus\footnote{\cite[alternative=entry][Gotzig2015Mar]}:

\startformula 
s=\frac{c_0\cdot t}{2}
 \stopformula

\startitemize[packed]
\item
  $s =$ Abstand
\item
  $c_0 =$ Geschwindigkeit der Energiewelle, bei Lidar:
  Lichtgeschwindigkeit (ca. $300 \frac{km}{s}$) und bei Radar:
  geschwindigkeit ausbreitung von elektromagnetischer Wellen (annähernd
  Lichtgeschwindigkeit)
\item
  $t =$ Zeit in $s$.
\stopitemize

\startformula 
\lambda = c \cdot T = \frac{c}{f}
 \stopformula


\chapter[title={Aufgabe 3},reference={aufgabe-3}]

\startframedtask
Ordnen Sie die in RADAR und LIDAR verwendeten Wellenlängen bzw.
Frequenzen in das elektromagnetische Spektrum ein.

https://www.leifiphysik.de/optik/elektromagnetisches-spektrum

\stopframedtask

\startuseMPgraphic{DiagonalRule}
  rulethickness := \frameddimension{rulethickness};

  drawoptions(
    withpen pencircle scaled rulethickness
    withcolor \MPcolor{\framedparameter{framecolor}});

  pair leftcorner, rightcorner;
  leftcorner  := (rulethickness, \overlayheight-rulethickness);
  rightcorner := (\overlaywidth-rulethickness, rulethickness);

  draw leftcorner -- rightcorner;
\stopuseMPgraphic

\defineoverlay
  [DiagonalRule]
  [\useMPgraphic{DiagonalRule}]

  \define[2]\DiagonalLabel{%
  \setuptabulate [after={\blank[\frameddimension{offset}]}]
  \starttabulate [|p|r|]
    \NC    \NC #2 \NC\NR
    \NC #1 \NC    \NC\NR
  \stoptabulate
}

\startplacetable[location=none]
\setupTABLE[r][1][style=bold,background=color,backgroundcolor=gray,width=2.5cm]
\setupTABLE[c][1][style=bold,background=color,backgroundcolor=gray]
\bTABLE
    \bTR
      \bTD [background=DiagonalRule]
        \DiagonalLabel{Kategorie}{Art}
      \eTD
      \bTD Radar \eTD
      \bTD Lidar \eTD
    \eTR
\startCSV
Wellenlänge:, $1m$, zwischen $1m$ und $780nm$
Frequenzen:, $300MHz$, $300GHz$ bis $385THz$
Art:, Radiowellen, Infrarot
\stopCSV
\eTABLE
\stopplacetable

\stopbodymatter

\startbackmatter
\component c_definitions
% \component c_bibliography
\stopbackmatter

\stoptext
